ch}) && \text{(canali)}\\
   \mid{}& \mathtt{par}\ c\ \mathtt{and}\ c\ [\mathtt{and}\ c]^{*}\ \mathtt{rap} && \text{(parallelo)}\\
   \mid{}& \mathtt{try}\ b\ c\ [\mathtt{rollback}\ c]\ \mathtt{yrt}\ b' && \text{(commit / rollback)}
\end{align*}

\noindent
Il bersaglio di un assegnamento è un \emph{luogo} $\ell ::= v \mid a[e]$: una
variabile intera oppure una cella di array. Le due forme si comportano allo
stesso modo — l'unica differenza è che la cella va individuata valutando
l'indice prima di scrivere.

\noindent
La sintassi concreta ammette due abbreviazioni. Una chiamata si può scrivere
senza parola chiave, $f(\bar v)$, e vale $\mathtt{call}\ f(\bar v)$. Il valore
$m$ di $\mathtt{local}$ e $\mathtt{delocal}$ può essere anche una variabile già
nello store, che si legge come in un'espressione: è la forma
$\mathtt{delocal}\ i = n$ che chiude un contatore arrivato al limite $n$.

\noindent
Un \emph{programma} è una lista di procedure
$\mathtt{procedure}\ f(\bar v)\ \{\, c \,\}$; $\bar v$ è la lista (eventualmente
vuota) dei parametri. Questa è la sintassi astratta: la sequenza si scrive qui
con il punto e virgola, $c_1\,;\,c_2$, come nelle regole di \S\ref{sec:seq},
mentre nel testo dei programmi i comandi si scrivono uno dopo l'altro, senza
separatori, e il corpo di una procedura arriva fino alla procedura successiva,
senza parentesi graffe (per convenzione un comando per riga, indentato). Il condizionale ha
\textbf{due} guardie ($b$ d'ingresso, $b'$ d'uscita) e il ciclo pure ($b_1$
d'ingresso, $b_2$ d'uscita): è ciò che rende il linguaggio reversibile
(\S\ref{sec:rev}). Le parentesi quadre segnano le clausole \emph{facoltative},
come in Janus: un ramo assente vale $\mathtt{skip}$, e nelle regole che seguono
si scrivono per esteso tutte e due, senza perdita di generalità.

\medskip
\noindent
L'esecuzione parte dalla procedura $\mathtt{main}$, senza parametri, dopo che
tutte le procedure del programma sono state registrate: il loro ordine nel file
non conta. A fine esecuzione l'interprete mostra le variabili di
$\mathtt{main}$. Un programma ben formato soddisfa inoltre questi vincoli,
che il compilatore verifica prima dell'esecuzione:
\begin{itemize}
  \item in un assegnamento $v \mathrel{\oplus}= e$ la variabile (o l'array) $v$
    non compare in $e$, né nell'indice della cella scritta
    (\S\ref{sec:assign}); $\mathtt{delocal}\ v = v$ non è ammesso;
  \item $\mathtt{local}$ e $\mathtt{delocal}$ formano un blocco, come in Janus:
    ogni $\mathtt{delocal}$ chiude il $\mathtt{local}$ aperto più di recente
    nella stessa sequenza di comandi, nessun $\mathtt{local}$ resta aperto alla
    fine della sequenza, e un $\mathtt{local}$ non ridichiara un parametro o un
    locale ancora aperto (\S\ref{sec:scope}); un $\mathtt{delocal}$ riporta il
    valore atteso e ha lo stesso tipo, e per un array la stessa lunghezza, del
    $\mathtt{local}$ che chiude;
  \item le dichiarazioni senza $\mathtt{local}$ ($\mathtt{int}\ x$,
    $\mathtt{stack}\ s$, $\mathtt{channel}\ c$) compaiono solo in
    $\mathtt{main}$: sono il suo store iniziale, a zero o vuoto, ed è ciò che
    l'interprete mostra a fine esecuzione; nelle altre procedure ogni variabile
    è un parametro o un locale;
  \item le procedure hanno nomi distinti, diversi da quelli delle operazioni
    predefinite ($\mathtt{show}$, $\mathtt{push}$, $\mathtt{pop}$,
    $\mathtt{ssend}$, $\mathtt{srecv}$) e da quelli che il compilatore riserva
    alle procedure che genera (che finiscono in \texttt{\_\_inv} o cominciano
    con \texttt{\_\_try\_}), e $\mathtt{main}$ non ha parametri;
  \item gli operandi hanno il tipo giusto, secondo la dichiarazione in scope
    (uno stesso nome può essere un intero in un blocco e uno stack in un
    altro): un'espressione legge interi e celle di array; un argomento di
    $\mathtt{call}$/$\mathtt{uncall}$ ha il tipo del parametro corrispondente;
    $\mathtt{push}$ e $\mathtt{pop}$ collegano una variabile intera e uno
    stack; $\mathtt{ssend}$ e $\mathtt{srecv}$ usano un canale e un messaggio
    non vuoto di variabilch}) && \text{(canali)}\\
   \mid{}& \mathtt{par}\ c\ \mathtt{and}\ c\ [\mathtt{and}\ c]^{*}\ \mathtt{rap} && \text{(parallelo)}\\
   \mid{}& \mathtt{try}\ b\ c\ [\mathtt{rollback}\ c]\ \mathtt{yrt}\ b' && \text{(commit / rollback)}
\end{align*}

\noindent
Il bersaglio di un assegnamento è un \emph{luogo} $\ell ::= v \mid a[e]$: una
variabile intera oppure una cella di array. Le due forme si comportano allo
stesso modo — l'unica differenza è che la cella va individuata valutando
l'indice prima di scrivere.

\noindent
La sintassi concreta ammette due abbreviazioni. Una chiamata si può scrivere
senza parola chiave, $f(\bar v)$, e vale $\mathtt{call}\ f(\bar v)$. Il valore
$m$ di $\mathtt{local}$ e $\mathtt{delocal}$ può essere anche una variabile già
nello store, che si legge come in un'espressione: è la forma
$\mathtt{delocal}\ i = n$ che chiude un contatore arrivato al limite $n$.

\noindent
Un \emph{programma} è una lista di procedure
$\mathtt{procedure}\ f(\bar v)\ \{\, c \,\}$; $\bar v$ è la lista (eventualmente
vuota) dei parametri. Questa è la sintassi astratta: la sequenza si scrive qui
con il punto e virgola, $c_1\,;\,c_2$, come nelle regole di \S\ref{sec:seq},
mentre nel testo dei programmi i comandi si scrivono uno dopo l'altro, senza
separatori, e il corpo di una procedura arriva fino alla procedura successiva,
senza parentesi graffe (per convenzione un comando per riga, indentato). Il condizionale ha
\textbf{due} guardie ($b$ d'ingresso, $b'$ d'uscita) e il ciclo pure ($b_1$
d'ingresso, $b_2$ d'uscita): è ciò che rende il linguaggio reversibile
(\S\ref{sec:rev}). Le parentesi quadre segnano le clausole \emph{facoltative},
come in Janus: un ramo assente vale $\mathtt{skip}$, e nelle regole che seguono
si scrivono per esteso tutte e due, senza perdita di generalità.

\medskip
\noindent
L'esecuzione parte dalla procedura $\mathtt{main}$, senza parametri, dopo che
tutte le procedure del programma sono state registrate: il loro ordine nel file
non conta. A fine esecuzione l'interprete mostra le variabili di
$\mathtt{main}$. Un programma ben formato soddisfa inoltre questi vincoli,
che il compilatore verifica prima dell'esecuzione:
\begin{itemize}
  \item in un assegnamento $v \mathrel{\oplus}= e$ la variabile (o l'array) $v$
    non compare in $e$, né nell'indice della cella scritta
    (\S\ref{sec:assign}); $\mathtt{delocal}\ v = v$ non è ammesso;
  \item $\mathtt{local}$ e $\mathtt{delocal}$ formano un blocco, come in Janus:
    ogni $\mathtt{delocal}$ chiude il $\mathtt{local}$ aperto più di recente
    nella stessa sequenza di comandi, nessun $\mathtt{local}$ resta aperto alla
    fine della sequenza, e un $\mathtt{local}$ non ridichiara un parametro o un
    locale ancora aperto (\S\ref{sec:scope}); un $\mathtt{delocal}$ riporta il
    valore atteso e ha lo stesso tipo, e per un array la stessa lunghezza, del
    $\mathtt{local}$ che chiude;
  \item le dichiarazioni senza $\mathtt{local}$ ($\mathtt{int}\ x$,
    $\mathtt{stack}\ s$, $\mathtt{channel}\ c$) compaiono solo in
    $\mathtt{main}$: sono il suo store iniziale, a zero o vuoto, ed è ciò che
    l'interprete mostra a fine esecuzione; nelle altre procedure ogni variabile
    è un parametro o un locale;
  \item le procedure hanno nomi distinti, diversi da quelli delle operazioni
    predefinite ($\mathtt{show}$, $\mathtt{push}$, $\mathtt{pop}$,
    $\mathtt{ssend}$, $\mathtt{srecv}$) e da quelli che il compilatore riserva
    alle procedure che genera (che finiscono in \texttt{\_\_inv} o cominciano
    con \texttt{\_\_try\_}), e $\mathtt{main}$ non ha parametri;
  \item gli operandi hanno il tipo giusto, secondo la dichiarazione in scope
    (uno stesso nome può essere un intero in un blocco e uno stack in un
    altro): un'espressione legge interi e celle di array; un argomento di
    $\mathtt{call}$/$\mathtt{uncall}$ ha il tipo del parametro corrispondente;
    $\mathtt{push}$ e $\mathtt{pop}$ collegano una variabile intera e uno
    stack; $\mathtt{ssend}$ e $\mathtt{srecv}$ usano un canale e un messaggio
    non vuoto di variabil